\section{Discussion}
\label{sec:discussion}

\subsection{What a ratchet without a control can and cannot tell you}

Both engines improved a validation score and both accepted on it. Neither
improvement survived a placebo. That gap has a mechanical explanation which
does not require either engine to be badly built.

An accept-if-strictly-better rule performs one implicit test per proposal,
keeps every winner, reports none of the tests, and has no way to ask whether a
difference could have arisen by chance. In our runs SkillOpt accepted its
winner on 21 of 21 validation items against a seed's 18, with no interval on
either. GEPA's return rule evaluates and selects on the full validation pool;
while our lineage file initially recorded the three-item reflection minibatch
the candidate was first proposed on, the search log committed with the run
(\Cref{sec:whatengines}) confirms that the winner was evaluated and selected
on the complete 21-item validation pool. A ratchet in that regime is a device
for selecting noise, and it will select noise reliably given enough proposals.

Neither engine is blind to overfitting by design. SkillOpt gates every edit
on a selection split disjoint from the rollouts the edit was generated from,
and GEPA maintains a validation pool for candidate selection. What neither
attaches is any notion of chance: each acceptance is a single
unreplicated evaluation, the same selection items are reused adaptively for
every decision in the run, and a strictly-better rule fires on a one-item
difference as readily as on a twenty-item one. Repeated evaluation at the
acceptance gate, or any test that asks how often a no-better candidate would
pass it, would close the gap, and nothing about either architecture prevents
that.

\subsection{Memorisation as the reachable optimum}

Both winners wrote training-item constants into the skill body, and both
landed on the same training template with different numbers
(\Cref{sec:whatengines}). Neither engine was asked to. Both were asked to
maximise a score on items they could see, and transcribing what they saw is
the best available way to do that.

This has a design consequence for anyone evaluating these engines. A holdout
that varies only the random seed of a template scores a transcribed template
\emph{rule} as generalisation, because the transcription still fits the
template it came from. A transcribed \emph{constant} does not survive a seed
change on this corpus, where every threshold is redrawn per item
(\Cref{sec:corpus}). \textbf{Split by template, or the rule channel stays
open.}

The constant ablation probe (\Cref{sec:whatengines}) demonstrates that the
transcription paid directly on the templates it targeted. Replacing SkillOpt's
hardcoded cutoffs for Loan Review with abstract notation dropped performance on
that template by 10.7 percentage points, below its matched placebo. A corpus
with fixed thresholds would have rewarded transcription outright. This one
redrew thresholds per seed, so the lifted numbers helped only on the template
instances that matched the search's training examples.

\subsection{The placebo and the seed skill}

Across both item sets, the placebo (matched on word count and structure but
carrying no decision procedure) scores at or above the initial seed skill. The
registration tests only the direction of an arm helping, so no formal $p$-value
attaches to that gap.

Read one way, this is a finding about the seed skill: on this corpus and model,
procedural instructions provided no measurable advantage over inert structured
text of the same length. Read another way, it is a methodological finding for
the wider literature: without a length-matched placebo control, any apparent
gain from an evolved skill is easily confounded with the mere presence of
structured text in the context window. Evaluating with dedicated placebos
disentangles instructional content from structural presence.

The largest effect this evaluation identified points to the harness question.
Nearly every unreadable answer from the seed skill and the placebos is a
generation that ran to the output cap without concluding; the seed skill
produces the most of any arm (\NUM{\atCapOn} across both sets, compared to
\NUM{\atCapPlacebo} for \texttt{placebo} and \NUM{\atCapOff} for \texttt{off};
\Cref{sec:parsefail}). On a single forced choice that is worth a few points and
reads as noise. In an agent loop, a model that will not stop is a different
kind of failure, and nothing here measures it.

\subsection{Where this leaves the availability and optimisation results}

Our result is closer to the near-zero prose effects of the controlled
SkillsBench-subset study~\citep{skillsbench2026} than to SkillOpt's $+23.5$
on GPT-5.5 in direct chat~\citep{skillopt2026}, and we report that with the
scope it deserves. We evaluated a 1.7B parameter local model across a
synthetic task family with computable ground truth. SkillOpt reports
comparable gains inside Codex and Claude Code, agent harnesses running
multi-step work, which is a different regime in a way that could matter: an
evolved skill that fixes a formatting pathology is
worth little on a single forced choice and could be worth a great deal across
a long trajectory. The runaway-generation result above is a concrete instance
of the same boundary.

\subsection{Why no larger-model replication exists}
\label{sec:noreplication}

The obvious next step is this study on a model large enough to carry a
verdict, and the corpus cannot host it. Screened with an empty prompt on
\NUM{\screenItems} items each, every hosted model in our committed screen
answers this corpus at or near ceiling: \NUM{\screenAtCeiling} of
\NUM{\screenModels} at \NUM{\screenBest} and the weakest at
\NUM{\screenWorst}, against \NUM{\accUnseenOff} for the study's local target.
At \NUM{\screenWorst} there is no headroom for any effect this design could
detect.

We then authored nine harder templates and none opened the headroom: on a 30B
model six hit the ceiling the originals hit, two turned out to be one-sided
in a way no arm could exploit, and the ninth had no room either. What the
attempt produced is an admission rule for future corpora and three
registered predictions falsified within a day of registration, which is the
pre-registration machinery working at authoring speed.
Appendix~\ref{sec:ceiling} reports the probes, their provenance, and the
rule.

\subsection{What a verdict is worth}

The claim we are entitled to is that on \NUM{\studyModel}, on this corpus,
under these controls, neither engine's output beat a document-shaped placebo
as registered. Harness variance exceeds model variance
by roughly $7.8\times$ on one factorial design~\citep{harnessdisclosure2026},
so transfer to another scaffold should not be assumed. That is the warning
that literature issues about everyone else's numbers, and we apply it to
ours.

The two checks this study contributes to that literature need no data at
all. Both are arithmetic over a registration, and either would have stopped
this study. A null from an underpowered design is not evidence of absence. A
null from a design with a $p$-floor above its own $\alpha$ is not evidence
of anything, and three of our six registered comparisons were exactly that.

\subsection{The checks, collected}

Every check this paper argues for runs on artifacts a study of this kind
already has.

\begin{enumerate}
  \item A placebo matched to the artifact under test, reported as the
    control (\Cref{sec:method}).
  \item An A/A repeat of one arm, scored into its own checkpoint
    (\Cref{sec:results}).
  \item The unit of analysis named at registration, and the $p$-floor
    $2^{-k}$ computed from the cluster count before the run
    (\Cref{sec:clustered}).
  \item The minimum detectable effect computed at the design's own
    clustering, read beside the effect the study later observes
    (\Cref{sec:limitations}).
  \item Discrimination reported beside accuracy, per template
    (\Cref{sec:signal}).
  \item The discordant split beside every paired $p$-value
    (\Cref{sec:method}).
  \item Every refused answer read once by a person, before a pattern in
    where the refusals fall is reported as a finding (\Cref{sec:parsefail}).
\end{enumerate}

\subsection{Negative results and boundary conditions}

The empirical findings bound what skill files and automated optimizers achieve
under controlled conditions. The seed skill scored at or below matched
placebos on both item sets. Template informedness screening revealed that one of
fourteen templates carried near-zero discrimination signal (\Cref{sec:signal}).
Two of eight pre-registered predictions did not reject the null
(\Cref{sec:predictions}). Refused response auditing eliminated spurious
non-null parse artifacts (\Cref{sec:parsefail}). Finally, ceiling screens on
frontier models demonstrated that single-prompt synthetic decision benchmarks
lack the headroom required to evaluate skill performance on larger models.
